% ============================================================
\chapter{Nonlinear Programming}
\label{ch:nonlinear}
% ============================================================

\section{Introduction}

When the objective function or constraints are nonlinear, we enter the realm
of \textbf{nonlinear optimization}. This chapter covers optimality conditions,
classical descent methods, and approaches for constrained problems.

\begin{intuition}
In linear programming, the optimum always lies at a vertex of the polyhedron.
In nonlinear optimization, the optimum can be anywhere --- inside the feasible
region, on its boundary, or even at a saddle point. More sophisticated
analytical tools are needed.
\end{intuition}

\section{Unconstrained optimization}

\subsection{Optimality conditions}

\begin{theorem}[First-order necessary condition]
\index{optimality conditions}
If $f : \R^n \to \R$ is differentiable and $x^*$ is a local minimum, then:
\[
  \nabla f(x^*) = 0.
\]
\end{theorem}

\begin{theorem}[Second-order sufficient condition]
If $f$ is twice differentiable, $\nabla f(x^*) = 0$, and the Hessian
$\nabla^2 f(x^*)$ is \textbf{positive definite}, then $x^*$ is a
\textbf{strict local minimum}.
\end{theorem}

\begin{warning}[title=\textbf{Local vs.\ global minima}]
For a nonconvex function, a point satisfying first- and second-order
conditions is only a \textbf{local} minimum. Only convexity of $f$ guarantees
that every local minimum is global.
\end{warning}

\subsection{Gradient descent}

\begin{definition}[Gradient descent]
\index{gradient descent}
\textbf{Gradient descent} generates a sequence $(x_k)$ by:
\[
  x_{k+1} = x_k - \alpha_k \nabla f(x_k),
\]
where $\alpha_k > 0$ is the \textbf{step size} (or learning rate).
\end{definition}

\begin{theorem}[Convergence for $L$-smooth convex functions]
If $f$ is convex and $\nabla f$ is $L$-Lipschitz, gradient descent with
constant step $\alpha = 1/L$ satisfies:
\[
  f(x_k) - f(x^*) \leq \frac{L \norm{x_0 - x^*}^2}{2k}.
\]
The convergence rate is $\mathcal{O}(1/k)$.
\end{theorem}

\begin{keyformulas}[title=\textbf{Step size selection}]
\begin{itemize}
  \item \textbf{Constant step}: $\alpha_k = \alpha < 2/L$ (simple,
        convergent).
  \item \textbf{Exact line search}:
        $\alpha_k = \argmin_{\alpha > 0} f(x_k - \alpha \nabla f(x_k))$.
  \item \textbf{Backtracking (Armijo)}: reduce $\alpha$ by factor
        $\beta \in (0,1)$ until
        $f(x_k - \alpha \nabla f(x_k)) \leq f(x_k) -
        c \alpha \norm{\nabla f(x_k)}^2$, $c \in (0,1)$.
\end{itemize}
\end{keyformulas}

\subsection{Newton's method}

\begin{definition}[Newton's method]
\index{Newton's method}
\textbf{Newton's method} uses second-order information:
\[
  x_{k+1} = x_k - [\nabla^2 f(x_k)]^{-1} \nabla f(x_k).
\]
\end{definition}

\begin{theorem}[Quadratic convergence]
If $f$ is twice continuously differentiable, $\nabla^2 f$ is Lipschitz, and
$x_0$ is sufficiently close to $x^*$ (where $\nabla^2 f(x^*)$ is positive
definite), the convergence is \textbf{quadratic}:
\[
  \norm{x_{k+1} - x^*} \leq C \norm{x_k - x^*}^2.
\]
\end{theorem}

\begin{remark}
Newton's method converges very fast near the solution but requires computing
and inverting the Hessian ($\mathcal{O}(n^3)$ per iteration).
\textbf{Quasi-Newton} methods (BFGS, L-BFGS) approximate the Hessian to
reduce cost.
\end{remark}

\begin{example}[Gradient descent vs.\ Newton]
Consider $f(x_1, x_2) = 10x_1^2 + x_2^2$ (ill-conditioned quadratic,
$L/\mu = 10$). From $x_0 = (10, 1)^\top$:
\begin{itemize}
  \item Gradient ($\alpha = 1/L = 1/20$): slow convergence, zigzag
        trajectory.
  \item Newton: converges in \textbf{one iteration} since $f$ is quadratic.
\end{itemize}
\end{example}

\section{Constrained optimization: KKT conditions}

\begin{definition}[Constrained optimization problem]
\[
  \min_{x \in \R^n} f(x) \quad \text{s.t.} \quad
  g_i(x) \leq 0,\; i = 1,\ldots,m, \qquad
  h_j(x) = 0,\; j = 1,\ldots,p.
\]
\end{definition}

\begin{theorem}[Karush-Kuhn-Tucker conditions]
\index{KKT conditions}
If $x^*$ is a local minimum and a constraint qualification holds (e.g., LICQ),
then there exist multipliers $\lambda^* \in \R^m_+$ and $\nu^* \in \R^p$
such that:
\begin{enumerate}
  \item \textbf{Stationarity}:
        $\nabla f(x^*) + \sum_{i} \lambda_i^* \nabla g_i(x^*)
        + \sum_{j} \nu_j^* \nabla h_j(x^*) = 0$.
  \item \textbf{Primal feasibility}: $g_i(x^*) \leq 0$, $h_j(x^*) = 0$.
  \item \textbf{Dual feasibility}: $\lambda_i^* \geq 0$.
  \item \textbf{Complementary slackness}: $\lambda_i^* g_i(x^*) = 0$.
\end{enumerate}
\end{theorem}

\begin{remark}
For convex problems (convex objective, convex inequality constraints, affine
equalities), the KKT conditions are also \textbf{sufficient} for global
optimality.
\end{remark}

\section{Penalty and barrier methods}

\subsection{Exterior penalty method}

\begin{definition}[Quadratic penalty]
\index{penalty method}
Replace the constrained problem by the sequence of unconstrained problems:
\[
  \min_{x}\; f(x) + \frac{\rho}{2} \sum_{i=1}^m [\max(0, g_i(x))]^2
  + \frac{\rho}{2} \sum_{j=1}^p [h_j(x)]^2,
\]
with $\rho \to +\infty$.
\end{definition}

\begin{theorem}[Penalty convergence]
If $\rho_k \to +\infty$ and $x_k$ is an (approximate) minimizer of the
penalized problem with parameter $\rho_k$, then every accumulation point of
$(x_k)$ solves the constrained problem.
\end{theorem}

\subsection{Barrier method (interior point)}

\begin{definition}[Logarithmic barrier]
\index{barrier method}
For inequality constraints $g_i(x) \leq 0$:
\[
  \min_{x}\; f(x) - \frac{1}{t} \sum_{i=1}^m \ln(-g_i(x)),
\]
with $t \to +\infty$. The barrier prevents leaving the interior of the
feasible region.
\end{definition}

\section{Sequential quadratic programming (SQP)}

\begin{definition}[SQP]
\index{SQP}
\textbf{Sequential Quadratic Programming} solves a nonlinear constrained
problem by solving at each iteration a \textbf{quadratic subproblem}:
\[
  \min_{d}\; \nabla f(x_k)^\top d + \frac{1}{2} d^\top H_k d \quad
  \text{s.t.} \quad g_i(x_k) + \nabla g_i(x_k)^\top d \leq 0,
\]
where $H_k$ approximates the Hessian of the Lagrangian. Then
$x_{k+1} = x_k + d_k$.
\end{definition}

\begin{remark}
SQP generalizes Newton's method to the constrained case. It is one of the most
efficient methods for medium-scale nonlinear optimization, with superlinear
convergence near the solution.
\end{remark}

\section{Python implementation}

\begin{codeblock}[title=\textbf{Nonlinear optimization with SciPy}]
\begin{minted}{python}
import numpy as np
from scipy.optimize import minimize

# Rosenbrock function
def rosenbrock(x):
    return (1 - x[0])**2 + 100*(x[1] - x[0]**2)**2

x0 = np.array([-1.0, 1.0])

# Quasi-Newton (L-BFGS-B)
res = minimize(rosenbrock, x0, method='L-BFGS-B')
print("L-BFGS-B:", res.x, "in", res.nit, "iterations")

# With constraints (SQP via SLSQP)
cons = ({'type': 'ineq', 'fun': lambda x: x[0] + x[1] - 1})
res_c = minimize(rosenbrock, x0, method='SLSQP',
                 constraints=cons)
print("SLSQP:", res_c.x)
\end{minted}
\end{codeblock}

\section{Exercises}

\begin{exercise}[$\star$ -- KKT conditions]
Find the KKT points of: $\min\; x_1^2 + x_2^2$ s.t.\ $x_1 + x_2 \geq 1$.
Verify that the solution is the global minimum.
\end{exercise}

\begin{exercise}[$\star\star$ -- Gradient descent]
Implement gradient descent with backtracking for
$f(x) = x_1^4 + x_2^4 - 2x_1 x_2$ from $x_0 = (2, 2)^\top$. Plot the
trajectory and the decrease of $f$.
\end{exercise}

\begin{exercise}[$\star\star$ -- Newton's method]
Apply Newton's method to $f(x) = e^{x_1 + x_2} + x_1^2 + x_2^2$ from
$x_0 = (1, 1)^\top$. Compute the first three iterations by hand.
\end{exercise}

\begin{exercise}[$\star\star\star$ -- Penalty method]
Solve $\min\; x_1^2 + x_2^2$ s.t.\ $x_1 + x_2 = 1$ by quadratic penalty
for $\rho \in \{1, 10, 100, 1000\}$. Observe convergence to the exact
solution.
\end{exercise}

\begin{exercise}[$\star\star\star$ -- SQP]
Formulate the QP subproblem of SQP for:
$\min\; (x_1 - 2)^2 + (x_2 - 1)^2$ s.t.\ $x_1^2 + x_2^2 \leq 1$.
Perform two iterations from $x_0 = (0, 0)^\top$.
\end{exercise}
